\documentclass[12pt]{article}
\usepackage[margin=1in,a4paper]{geometry}
\usepackage{amsmath,amssymb,fullpage,xcolor,mhchem,mdframed}
\DeclareMathOperator\Tr{Tr}
\begin{document}
\title{Chapter 41: Feynman Rules for Dirac Fields}
\date{\today}
\maketitle

\begin{mdframed}[backgroundcolor=blue!5]
\textbf{Note:} This chapter applies Ch.37's canonical quantization to Dirac
fields. The Feynman propagator and spin sums here are used in every fermion-loop
and external-leg calculation throughout the rest of the course.
\end{mdframed}

\section{Dirac Propagator}\label{sec:propagator}

The Feynman propagator is the time-ordered vev:
\begin{equation}
S_F(x-y) = \langle0| T\{\psi(x)\bar\psi(y)\}|0\rangle .
\end{equation}
Evaluating via the mode expansion gives, in momentum space:
\begin{mdframed}[backgroundcolor=green!5]
\begin{equation}
S_F(k) = \frac{i(k\!\!\!/+m)}{k^2-m^2+i\varepsilon}
      = \frac{i}{k\!\!\!/-m+i\varepsilon} .
\tag{39.3}
\end{equation}
\end{mdframed}

The Feynman $i\varepsilon$ prescription places poles at $k^0 = \pm E_\mathbf{k} \mp i\varepsilon$,
specifying which contour deformation gives the correct time ordering.

\section{External Fermion Wave Functions}\label{sec:external}

For a fermion with momentum $p^\mu=(E,\mathbf{p})$ and spin label $s=\pm$:
\begin{equation}
u_s(p) = \begin{pmatrix}
\sqrt{E+\mathbf{p}\cdot\boldsymbol{\sigma}}\,\chi_s \\
\sqrt{E-\mathbf{p}\cdot\boldsymbol{\sigma}}\,\chi_s
\end{pmatrix},
\qquad
\chi_+ = \begin{pmatrix}1\\0\end{pmatrix},\;
\chi_- = \begin{pmatrix}0\\1\end{pmatrix} .
\tag{39.5}
\end{equation}
In the massless limit and for $\mathbf{p}=(0,0,E)$ (along $+z$):
\begin{equation}
u_+(p) = \sqrt{2E}\begin{pmatrix}1\\0\\1\\0\end{pmatrix},
\qquad
u_-(p) = \sqrt{2E}\begin{pmatrix}0\\1\\0\\-1\end{pmatrix} .
\tag{39.6}
\end{equation}

\section{Spin Sums}\label{sec:spin_sums}

The completeness (spin-sum) identities are:
\begin{mdframed}[backgroundcolor=yellow!5]
\begin{align}
\sum_s u_s(p)\,\bar u_s(p) &= p\!\!\!/+m , \tag{39.8}\\
\sum_s v_s(p)\,\bar v_s(p) &= p\!\!\!/-m . \tag{39.9}
\end{align}
\end{mdframed}

For unpolarized cross sections, each initial fermion spin contributes a factor $\frac12$.

\section{Trace Theorems}\label{sec:traces}

The fundamental trace identities (4D Minkowski, $\eta=\mathrm{diag}(-1,+1,+1,+1)$):
\begin{mdframed}[backgroundcolor=red!5]
\begin{align}
\Tr[\mathbf{1}] &= 4 , \tag{39.10}\\
\Tr[\gamma^\mu\gamma^\nu] &= 4\eta^{\mu\nu} , \tag{39.12}\\
\Tr[\gamma^\mu\gamma^\nu\gamma^\rho\gamma^\sigma]
&= 4\bigl(\eta^{\mu\nu}\eta^{\rho\sigma}
          -\eta^{\mu\rho}\eta^{\nu\sigma}
          +\eta^{\mu\sigma}\eta^{\nu\rho}\bigr) . \tag{39.14}
\end{align}
\end{mdframed}

\begin{center}                                                                                                                                                                                              
\textbf{Numpy verification of trace theorems:} 
\textbf{Numpy verification of trace theorems:}
\end{center}
\begin{center}
\begin{tabular}{|l|c|}
\hline
Check & Result \\
\hline
$\Tr[\gamma^\mu\gamma^\nu] = 4\eta^{\mu\nu}$ & FAIL \\
$\Tr[\gamma^\mu\gamma^\nu\gamma^\rho\gamma^\sigma]$ identity & PASS \\
$\Tr[\gamma^5\gamma^\mu\gamma^\nu\gamma^\rho\gamma^\sigma] = -4i\varepsilon^{\mu\nu\rho\sigma}$ & See Ch.41 \\
\hline
\end{tabular}
\end{center}

\section{QED Feynman Rules}\label{sec:qed_rules}

\begin{center}
\begin{tabular}{|l|c|}
\hline
Object & Rule \\
\hline
Photon propagator & $\displaystyle\frac{-i\eta_{\mu\nu}}{q^2+i\varepsilon}$ \\
Fermion propagator & $\displaystyle\frac{i(p\!\!\!/+m)}{p^2-m^2+i\varepsilon}$ \\
QED vertex & $-ie\,\gamma_\mu$ \\
External fermion & $u_s(p)$, $\bar u_s(p)$ \\
External antifermion & $v_s(p)$, $\bar v_s(p)$ \\
External photon & $\varepsilon_\mu(k,\lambda)$ \\
\hline
\end{tabular}
\end{center}

\section{Spinor Completeness (Massless)}\label{sec:spin_completeness}

For massless $p^\mu=(E,0,0,E)$ along $+z$, the spin-sum gives:
\begin{mdframed}
\begin{align*}
\sum_s u_s\bar u_s
&= p\!\!\!/ = \begin{pmatrix}
0&0&E&0\\0&0&0&-E\\E&0&0&0\\0&-E&0&0
\end{pmatrix} .
\end{align*}
Numpy difference from $p\!\!\!/$: $\|$diff$\|_F = $4.90e+00 \\
(Zero = PASS)
\end{mdframed}

\section{Significance for the MHV Program}

Ch.39 is foundational because:
\begin{enumerate}
\item Every QCD amplitude with external quarks uses the spin-sum formula
  $\sum u\bar u = p\!\!\!/+m$.
\item The trace theorems reduce loop calculations to scalar products of momenta.
\item In the massless chiral limit, the Dirac spinor basis becomes the
  helicity-spinor basis of Ch.42: $u_+(k)\to\lambda$, $u_-(k)\to\tilde\lambda$.
\end{enumerate}

\end{document}
